\subsection{Phân tích các bên liên quan}

Việc phát triển và triển khai Hệ thống Thương mại Điện tử cho Bán lẻ Điện thoại tích hợp Trợ lý Ảo AI Agent là một bước tiến công nghệ quan trọng, mang lại những lợi ích thiết thực và giá trị gia tăng rõ rệt cho các bên liên quan trực tiếp đến hoạt động kinh doanh và vận hành hệ thống. 

\subsubsection{Khách hàng}

Đối với người dùng cuối, hệ thống cung cấp một nền tảng mua sắm trực tuyến hiện đại và thuận tiện. Lợi ích chính nằm ở việc cá nhân hóa trải nghiệm thông qua trợ lý AI Agent, giúp khách hàng nhanh chóng tìm được sản phẩm phù hợp với nhu cầu. Hệ thống hỗ trợ giải đáp thắc mắc tức thì, đơn giản hóa quy trình đặt hàng và đảm bảo tính minh bạch trong việc theo dõi trạng thái đơn hàng cũng như bảo mật thông tin cá nhân.

\subsubsection{Đội ngũ vận hàng và quản trị viên}

Đối với đội ngũ nội bộ, hệ thống mới giúp tăng hiệu suất làm việc và tối ưu hóa quy trình quản lý. Lợi ích chính là việc tự động hóa khâu chăm sóc khách hàng, giúp đội ngũ hỗ trợ giảm tải khối lượng câu hỏi thường gặp. Quản trị viên có được một giao diện tập trung để quản lý sản phẩm, giá cả và khuyến mãi hiệu quả. Hơn nữa, hệ thống cung cấp báo cáo cơ bản và cơ chế cập nhật trạng thái đơn hàng, hỗ trợ công tác quản lý và ra quyết định.

\subsubsection{Chủ sở hữu doanh nghiệp}

Về mặt kinh doanh, hệ thống giúp tối ưu hóa khả năng tiếp cận khách hàng trên nền tảng số và nâng cao hiệu suất bán hàng. Việc ứng dụng công nghệ AI Agent góp phần tạo ra sự khác biệt trong dịch vụ, hỗ trợ tăng tỷ lệ chốt đơn và xây dựng lòng trung thành của khách hàng. Về mặt kỹ thuật, kiến trúc vi dịch vụ giúp hệ thống vận hành ổn định, dễ dàng nâng cấp hoặc mở rộng các tính năng mới theo nhu cầu thực tế mà không gây gián đoạn các dịch vụ sẵn có.

\subsubsection{Nhóm phát triển đề tài}

Dự án là nền tảng để nhóm phát triển ứng dụng và kiểm chứng các kỹ thuật hiện đại, đặc biệt là trong lĩnh vực Xử lý Ngôn ngữ Tự nhiên (NLP) và thiết kế kiến trúc vi dịch vụ cho thương mại điện tử. Lợi ích là kinh nghiệm thực tiễn trong việc tích hợp thành công công nghệ AI vào một hệ thống giao dịch phức tạp, tạo ra một sản phẩm minh chứng có giá trị nghiên cứu và ứng dụng cao.

Từ việc phân tích các bên liên quan, có thể thấy nhu cầu của từng nhóm được cụ thể hóa rõ ràng và gắn liền với vai trò mà họ đảm nhận trong hệ thống.Những nhu cầu trên chính là cơ sở để xác định các yêu cầu chức năng của hệ thống, bao gồm tính năng tìm kiếm, giỏ hàng, thanh toán mô phỏng, trợ lý ảo kịch bản và quản trị sản phẩm. Song song, các yêu cầu phi chức năng ở mức tối thiểu như sự ổn định cho bản demo, tốc độ phản hồi hợp lý và bảo mật thông tin người dùng cũng cần được đảm bảo để hệ thống có thể vận hành hiệu quả trong phạm vi nghiên cứu.